\section{Заключение}
\label{sec:Section5} \index{Section5}

В данной работе разработана система автоматического анализа и верификации правил фильтрации сетевого трафика, выявляющая ошибки конфигурации межсетевых экранов. Система реализована в виде прототипа \textit{RuleAnalyzer} и интегрирована в систему инвентаризации и валидации сетевых инфраструктур.

Были выполнены следующие задачи:
\begin{enumerate}
    \item Проведено исследование существующих методов анализа правил фильтрации; выявлено, что известные подходы рассчитаны под конкретную платформу или ограниченный набор полей.
    \item Разработана унифицированная вендор-независимая модель представления правил фильтрации и объектов политики безопасности, охватывающая как классические списки доступа, так и политики межсетевых экранов нового поколения, включая отечественные решения.
    \item Выбран и адаптирован для анализа метод формальной верификации FIREMAN.
    \item Реализованы алгоритмы обнаружения ошибок конфигурации. Помимо классических аномалий между правилами --- затенения, избыточности, корреляции и обобщения --- набор обнаруживаемых отклонений расширен: добавлены аномалии уровня одного правила (выключенные, истёкшие и неиспользуемые правила) и ошибки уровня объектов политики безопасности. Для правил с невыразимыми в модели квалификаторами применён консервативный анализ через разрешающее и запрещающее замыкания, не порождающий ложных критических ошибок.
    \item Разработан прототип системы анализа правил фильтрации и выполнена его интеграция в систему инвентаризации и валидации сетевых инфраструктур, разрабатываемую в ООО <<Дипси Нетворкс>>.
    \item Проведено нагрузочное тестирование, показавшее применимость прототипа к анализу больших наборов правил межсетевых экранов нового поколения.
\end{enumerate}

Планируемые будущие исследования по данной теме:

\textbf{Инкрементальный анализ.} В текущем виде анализатор проверяет конфигурацию, уже развёрнутую на устройстве. Более ценным сценарием является проверка \emph{до} развёртывания (\emph{pre-change analysis}), когда администратор редактирует набор правил и хочет убедиться в отсутствии ошибок перед его применением. Пересчитывать анализ всего набора на каждое изменение нерационально, поэтому необходимо по внесённому изменению определять затронутую им часть набора и пересчитывать только её, сохраняя остальной результат. Это позволило бы встроить проверку прямо в процесс редактирования конфигурации и обнаруживать ошибки на этапе подготовки изменений, а не их эксплуатации.

\textbf{Анализ достижимости.} Разработанный анализатор проверяет \emph{внутреннюю} согласованность набора правил одного устройства. Другой класс задач --- анализ достижимости (см.~\ref{sec:Section2:reachability}): определение того, какой трафик может пройти от одного узла сети до другого. В отличие от поиска аномалий, эта задача рассматривает не отдельное устройство, а сеть целиком --- последовательность межсетевых экранов на пути трафика с учётом топологии и маршрутизации. Унифицированная модель и аппарат предикатов над пространством заголовков, построенные в работе, применимы и здесь, однако переход к анализу достижимости требует моделирования трансляции адресов (NAT): она изменяет заголовки пакетов между устройствами и в настоящей работе сознательно оставлена за рамками. Развитие анализатора в эту сторону позволило бы перейти от выявления ошибок в наборе правил к проверке сквозной политики доступа в сети.

\newpage
